\chapter{Fazit}
\label{chap:fazit}
\section{Zusammenfassung}

Ziel dieser Arbeit war es, eine framework-unabhängige Abstraktionsschicht für die
Definition parametrisierter Quantenschaltkreise und Observablen zu entwerfen, im
\textit{Executor} zu implementieren und ihren Laufzeit-Overhead gegenüber der
direkten Nutzung der Frameworks zu bestimmen. Dieses Ziel wurde erreicht.

Kapitel~\ref{chap:umsetzung} beantwortet FF1. Die Abstraktion ruht auf vier
framework-unabhängigen Datentypen, deren zulässiges Vokabular zentral in
\texttt{SUPPORTED\_GATES} festgelegt und durch einen Paritätstest über beide
Backends abgesichert ist. Die bisherige Bindung an \textit{Qiskit} wurde damit
aufgelöst, da weder die Basisklassen noch das \textit{PennyLane}-Backend
\textit{Qiskit} importieren.

Kapitel~\ref{chap:vorteile} beantwortet FF2. Als qualitative Vorteile wurden eine
einheitliche Definitionsschnittstelle, die Modularisierbarkeit von Schaltkreisen,
optionale Framework-Abhängigkeiten sowie die Entkopplung von Framework-Versionen
herausgearbeitet. Zugleich wurde gezeigt, dass sich die Ausdrucksstärke der
Schnittstelle notwendigerweise auf die Schnittmenge der unterstützten Frameworks
beschränkt.

Kapitel~\ref{chap:benchmark} beantwortet FF3 mit einem zweigeteilten Befund.
Gegenüber \textit{PennyLane} entsteht in der Ausführung kein messbarer Overhead, da
der \textit{Executor} eine vorbereitete Instruktionsliste abspielt, während die
native Seite den Schaltkreis bei jedem Aufruf neu aufzeichnet. Gegenüber
\textit{Qiskit} kostet die Abstraktion dagegen messbar, und zwar um den Faktor
$1{,}97$ bei der Erstellung, $1{,}40$ bei der Ausführung und $2{,}40$ bei der
Gradientenberechnung. Der relative Overhead der Ausführung nimmt mit wachsender
Qubit-Zahl ab, da der Zusatzaufwand aus reiner Python-Buchhaltung besteht und nicht
von der Registergröße abhängt. Allein bei der Gradientenberechnung über
\textit{Qiskit} wächst der Faktor mit der Gatterzahl, was auf die tiefen
Schaltkreiskopien der Parameter-Shift-Regel zurückgeht. Der gemessene Vorsprung des
\textit{Executors} bei der Gradientenberechnung über \textit{PennyLane} bleibt wegen
abweichender Bibliotheksversionen zwischen beiden Messläufen nur eingeschränkt
belastbar.

Der gemessene Mehraufwand fällt gegenüber dem gewonnenen Nutzen gering aus.
Gegenüber \textit{PennyLane} arbeitet der \textit{Executor} auf dem Niveau der
nativen API, und gegenüber \textit{Qiskit} bewegt sich der Overhead in der
Ausführung in einem Bereich, der sich mit wachsender Registergröße weiter
angleicht, da der Zusatzaufwand nicht mit der Qubit-Zahl skaliert. Er beschränkt
sich damit auf einen konstanten Anteil an Python-Arbeit, dem eine einmalige, über
sämtliche Backends hinweg wiederverwendbare Problembeschreibung gegenübersteht.
Gerade dort, wo Laufzeit tatsächlich ins Gewicht fällt, also bei großen Registern
und rechenintensiven Simulationen, verliert der Overhead zusätzlich an Bedeutung.
Die in Kapitel~\ref{chap:vorteile} beschriebenen Vorteile, insbesondere die
kontrollierte Vergleichbarkeit von Backends, die optionalen
Framework-Abhängigkeiten und die Stabilität über Versionswechsel hinweg, wiegen
diesen Anteil in der Praxis auf.

Damit sind alle drei Forschungsfragen beantwortet.
% ----------------------------------------------------------------
\newpage
\section{Ausblick}
\label{sec:ausblick}
Der Executor im Allgemeinen und die Abstraktionsschicht können noch
weiterentwickelt werden, um Anbindungen an weitere Backends zu schaffen. Im
Moment werden nur \textit{PennyLane} und \textit{Qiskit} als API-Endpunkte
unterstützt, was bereits eine gute Basis bildet. Darüber hinaus ist die Grundlage
für eine Implementierung von \textit{Qulacs} und der Pauli-Propagation bereits
gegeben. Bei diesen fehlt jedoch noch die Umwandlung der abstrakten Schaltkreis-
und Observablen-Objekte in die Framework-spezifischen Repräsentationen. Auch
weitere Frameworks wie \textit{Cirq}, \textit{Braket SDK} oder \textit{Q\#} wären
sinnvolle Ergänzungen. Damit würde der Executor ein größeres Ökosystem
unterstützen und mehr Nutzerinnen und Nutzer ansprechen, die diese Backends
einsetzen.

Zudem orientiert sich die Abstraktionsschicht derzeit an der NISQ-Ära, also an verrauschten Quantencomputern ohne Fehlerkorrektur~\cite{preskill2018quantum}. Ihr Entwurf ist daher auf parametrisierte Schaltkreise sowie die Berechnung von Erwartungswerten und Gradienten ausgelegt. Perspektivisch wird jedoch das fehlertolerante Quantencomputing wichtiger, das mit fehlerkorrigierten logischen Qubits arbeitet. Diese Technologie steht allerdings noch am Anfang und ist bisher nicht praktisch nutzbar, sodass eine Anpassung der Abstraktionsschicht erst mit der Verfügbarkeit entsprechender Backends sinnvoll wird.

% Was sind sinnvolle nächste Schritte?
% Offen gebliebene Fragen, mögliche Erweiterungen,
% Verbesserungspotenziale, Anknüpfungspunkte für weitere Arbeiten.